% The butterfly of Section~\ref{sec:polar} at $N = 4$, drawn by hand. Input by
% textbook.tex; a sketch and not a rendered figure, so it is outside the figure
% cache and the render cap, which govern generated output only.
\begin{figure}[htbp]
\centering
\begin{tikzpicture}[
    scale=1.0,
    source/.style={fgvariable, minimum size=0.5cm, font=\scriptsize},
    kernel/.style={fgdata, minimum size=0.42cm, inner sep=1pt,
                   font=\scriptsize},
    frozen/.style={fgvariable, fill=black!12, minimum size=0.5cm,
                   font=\scriptsize}
]
  % the source vector: two frozen, two carrying the message
  \node[frozen] (u0) at (0, 2.4) {$u_0$};
  \node[frozen] (u1) at (0, 1.6) {$u_1$};
  \node[source] (u2) at (0, 0.8) {$u_2$};
  \node[source] (u3) at (0, 0.0) {$u_3$};

  % two stages of the kernel, one column each
  \foreach \i/\y in {0/2.4, 1/1.6, 2/0.8, 3/0.0} {
    \node[kernel] (a\i) at (1.7, \y) {$\oplus$};
    \node[kernel] (b\i) at (3.4, \y) {$\oplus$};
    \draw[fglink] (u\i) -- (a\i);
    \draw[fglink] (a\i) -- (b\i);
  }
  % the kernel pairs each column: stage one within halves, stage two across
  \draw[fglink] (u0) -- (a1);
  \draw[fglink] (u2) -- (a3);
  \draw[fglink] (a0) -- (b2);
  \draw[fglink] (a1) -- (b3);

  % the transmitted bits and their channel scores
  \foreach \i/\y in {0/2.4, 1/1.6, 2/0.8, 3/0.0} {
    \node[source] (c\i) at (5.1, \y) {$c_{\i}$};
    \node[fgdata, minimum size=0.42cm, font=\scriptsize] (l\i) at (6.4, \y)
      {$L_{\i}$};
    \draw[fglink] (b\i) -- (c\i);
    \draw[fglink] (c\i) -- (l\i);
  }

  \node[fgnote] at (0, -0.7) {source};
  \node[fgnote] at (2.55, -0.7) {$n = 2$ stages of $\oplus$};
  \node[fgnote] at (5.75, -0.7) {channel};
\end{tikzpicture}
\caption{The transform of Equation~\ref{eq:polar-transform} at $N = 4$ as a
factor graph. Each $\oplus$ is one application of the kernel: a parity factor
over the pair it joins. Shaded source variables are frozen to zero and carry
no message, and which they are is what Equation~\ref{eq:polar-bec} chooses.
Successive cancellation is a schedule on this graph and not a separate
algorithm --- it decides $u_0$, then $u_1$ given that decision, and so on,
each pass exact given the ones before it.}
\label{fig:polar:sketch}
\end{figure}
